\section*{Acknowledgments}\label{sec:ack}

This paper is dedicated to nature and all her beauty.

The rose has served across traditions as symbol of the
Divine Mother---Tripurasundari in the \textit{Sri Vidya},
the Gothic rosette, the Sufi \textit{gul}---representing the
manifestation of perfection in material form.
The structure is $\golden^2 = \golden + 1$: the rose's five
petals trace the pentagon, her spirals follow the Fibonacci
sequence $F_n/F_{n-1} \to \golden$, and each petal is placed at
the golden angle $360^\circ/\golden^2 \approx 137.5^\circ$.
In the language of Hegel's \textit{Encyclop\"adie}~(\S247),
nature is the Idea in the form of otherness---the rational
structure externalised and asleep in matter.  The mathematics
that arises from observing her is the return: the Idea
awakening to itself.  This paper traces that return from the
pentagon to $\zeta(2k+1)$.

As a companion paper to~\cite{V11}, this work shares the
foundational, institutional, and technical acknowledgments
therein.
